\documentclass[11pt]{article}
\usepackage{mathblatt}
% gesetzt von werkzeuge/setzer.py; Lösungen nach bau/bauregeln.md „Lösungen“.
\newcommand{\kennung}{HDK}
\begin{document}
\pfheftstil
\fancyfoot[C]{{\footnotesize\color{mbgrau}\kennung\hfill\thepage}}
\pfheftkopf{Winkelsummen, Dreiecke und Vierecke}{Klasse 7 \textperiodcentered{} Lösungen}

\begin{pfloesung}{}
\lz{1b)}{$\gamma = 47^\circ$}{$\gamma = 180^\circ - 38^\circ - 95^\circ = 47^\circ$}{}
\end{pfloesung}
\begin{pfloesung}{}
\lz{2.}{$\beta = 66{,}7^\circ$}{$\beta = 180^\circ - 63{,}5^\circ - 49{,}8^\circ = 66{,}7^\circ$}{}
\end{pfloesung}
\begin{pfloesung}{}
\lz{3.}{$\gamma = 62{,}5^\circ$}{$\gamma = 180^\circ - 90^\circ - 27{,}5^\circ = 62{,}5^\circ$}{}
\end{pfloesung}
\begin{pfloesung}{}
\lz{4.}{Nein: die dritte Ecke hat nur $54^\circ$.}{Nein; dritte Ecke $180^\circ - 66{,}5^\circ - 59{,}5^\circ = 54^\circ$; $54^\circ < 55^\circ$: zu spitz.}{}
\end{pfloesung}
\begin{pfloesung}{}
\lz{5b)}{$\alpha = 69^\circ$, $\gamma = 42^\circ$}{$\alpha = \beta = 69^\circ$ (Basiswinkel); $\gamma = 180^\circ - 69^\circ - 69^\circ = 42^\circ$}{}
\end{pfloesung}
\begin{pfloesung}{}
\lz{6.}{$\beta = \gamma = 71^\circ$}{$\beta = \gamma = (180^\circ - 38^\circ) : 2 = 71^\circ$}{}
\end{pfloesung}
\begin{pfloesung}{}
\lz{7.}{$BC = 5$\,cm, $\gamma = 50^\circ$}{$BC = 5$\,cm (gleiche Winkel bei $A$ und $B$, gegenüber liegen $BC$ und $AC$); $\gamma = 180^\circ - 65^\circ - 65^\circ = 50^\circ$}{}
\end{pfloesung}
\begin{pfloesung}{}
\lz{8.}{Sicher: $74^\circ$ an jedem Holm.}{Holme und Boden bilden ein gleichschenkliges Dreieck mit der Spitze oben: $(180^\circ - 32^\circ) : 2 = 74^\circ$; $74^\circ > 70^\circ$.}{}
\end{pfloesung}
\begin{pfloesung}{}
\lz{9b)}{$\gamma = 61^\circ$}{$\gamma = 360^\circ - 78^\circ - 125^\circ - 96^\circ = 61^\circ$}{}
\end{pfloesung}
\begin{pfloesung}{}
\lz{10.}{$\beta = \delta = 95^\circ$}{$\beta = \delta = (360^\circ - 52^\circ - 118^\circ) : 2 = 95^\circ$}{}
\end{pfloesung}
\begin{pfloesung}{}
\lz{11.}{Alle vier Seiten sind gleich lang.}{(Raute: Parallelogramm mit vier gleich langen Seiten; gleiche Winkel hat nur das Quadrat.)}{}
\end{pfloesung}
\begin{pfloesung}{}
\lz{12.}{Nein: links und rechts je $108^\circ$ (nicht $106^\circ$).}{Nein; $84^\circ + 60^\circ + 106^\circ + 106^\circ = 356^\circ$, nicht $360^\circ$. Richtig: links und rechts je $(360^\circ - 84^\circ - 60^\circ) : 2 = 108^\circ$.}{}
\end{pfloesung}
\begin{pfloesung}{}
\lz{13b)}{$\varepsilon = 27^\circ$, $\varphi = 42^\circ$, $\gamma = 69^\circ$}{$\varepsilon = 180^\circ - 90^\circ - 63^\circ = 27^\circ$; $\varphi = 180^\circ - 90^\circ - 48^\circ = 42^\circ$; $\gamma = 27^\circ + 42^\circ = 69^\circ$}{}
\end{pfloesung}
\begin{pfloesung}{}
\lz{14.}{$\varepsilon = 26^\circ$}{Winkel $ADC = 180^\circ - 64^\circ = 116^\circ$ (Nebenwinkel); $\varepsilon = 180^\circ - 38^\circ - 116^\circ = 26^\circ$}{}
\end{pfloesung}
\begin{pfloesung}{}
\lz{15.}{zweimal $45^\circ$, zusammen $90^\circ$}{Dreieck $ADC$ ist gleichschenklig ($AD = DC$) mit $90^\circ$ bei $D$: Basiswinkel $(180^\circ - 90^\circ) : 2 = 45^\circ$, also Winkel $ACD = 45^\circ$. Ebenso im Dreieck $DBC$: Winkel $DCB = 45^\circ$. $\gamma = 45^\circ + 45^\circ = 90^\circ$.}{}
\end{pfloesung}
\begin{pfloesung}{}
\lz{16.}{Nur Zelt 1: oben $45^\circ + 45^\circ = 90^\circ$; bei Zelt 2 rechts nicht $45^\circ$.}{Zelt 1: Stange und Boden bilden links und rechts je ein rechtwinkliges Dreieck mit zwei gleich langen Seiten (1,5\,m). Seine Basiswinkel sind $(180^\circ - 90^\circ) : 2 = 45^\circ$, oben also $45^\circ + 45^\circ = 90^\circ$: ein rechter Winkel. Zelt 2: Links ist es wie bei Zelt 1, oben $45^\circ$. Rechts sind die Seiten 1,5\,m und 2\,m nicht gleich lang, also sind die beiden Winkel dort nicht gleich und oben nicht $45^\circ$. Zusammen nicht $90^\circ$: kein rechter Winkel.}{}
\end{pfloesung}
\begin{pfloesung}{}
\lz{17.}{$61{,}7^\circ$}{$180^\circ - 46{,}5^\circ - 71{,}8^\circ = 61{,}7^\circ$}{}
\end{pfloesung}
\begin{pfloesung}{}
\lz{18.}{$\delta = 109^\circ$}{$\delta = 360^\circ - 71^\circ - 64^\circ - 116^\circ = 109^\circ$}{}
\end{pfloesung}
\begin{pfloesung}{}
\lz{19.}{$\gamma = 71^\circ$}{$\beta = 180^\circ - 118^\circ = 62^\circ$ (Nebenwinkel); $\gamma = 180^\circ - 47^\circ - 62^\circ = 71^\circ$}{}
\end{pfloesung}
\begin{pfloesung}{}
\lz{20.}{Die Diagonalen stehen senkrecht aufeinander.}{}{}
\end{pfloesung}
\begin{pfloesung}{}
\lz{21.}{Haus 2, um $3^\circ$ ($38^\circ$ gegen $35^\circ$).}{Haus 2: Neigung $(180^\circ - 104^\circ) : 2 = 38^\circ$; $38^\circ - 35^\circ = 3^\circ$: Das Dach von Haus 2 ist um $3^\circ$ steiler.}{}
\end{pfloesung}
\end{document}
